\subsubsection*{Lines in the Cartesian plane}

A line is born from one fixed point and one allowed direction:
\[
L=\{P_0+t\mathbf v:t\in\mathbb R\}.
\]
The parameter $t$ records how far and in which orientation we move along the
permitted direction. The normal equation
\[
\mathbf n\cdot(P-P_0)=0
\]
arises from an equivalent demand: every allowed displacement must have zero
component in a perpendicular direction. The slope form $y=mx+b$ then appears as
a useful but limited coordinate consequence, not as the universal meaning of a
line.

The section teaches the student to choose a representation according to the
question. Parametric form describes motion. Normal form expresses membership and
perpendicularity. Slope form makes certain coordinate features easy to read.
Intersections with axes, the line through two points, and slope criteria become
consequences of the vector construction rather than unrelated rules.

Projection provides a further application. The nearest point on a line is obtained
by retaining the component parallel to the line and removing the perpendicular
remainder. The point-to-line distance formula is therefore the length of a normal
component. The student should be able to derive and move among the principal
forms of a line, select the useful form for a problem, test parallelism and
perpendicularity, find intersections, project a point onto a line, and justify
the distance formula.
